\documentclass[tikz,border=8pt]{standalone}
\usepackage{ctex}
\usepackage{amsmath}
\usetikzlibrary{calc, arrows.meta}

\begin{document}
% 正三棱锥 P-ABC，E 是 BC 中点，PE ⊥ 平面 ABC
% 使用等距投影（oblique/cabinet）
\begin{tikzpicture}[
  thick,
  dot/.style={fill=black, circle, inner sep=1.5pt},
]

% 正三棱锥的底面 ABC（等边三角形，俯视投影）
% A 在左下，B 在右下，C 在正上（投影后）
% 投影：等腰三角形底部，A(-1.5,-0.3), B(1.5,-0.3), C(0,1.5)
\coordinate (A) at (-1.8, 0);
\coordinate (B) at (1.8, 0);
\coordinate (C) at (0, 1.4);
\coordinate (P) at (0, 4.2);   % 顶点
\coordinate (E) at (0.9, 0.7); % BC 中点

% 画底面（浅蓝色填充）
\fill[blue!10, draw=blue!40, fill opacity=0.7] (A) -- (B) -- (C) -- cycle;

% 底面三条边
\draw[blue!60] (A) -- (B) node[midway, below, font=\small] {};
\draw[blue!60] (B) -- (C);
\draw[blue!60] (A) -- (C);

% 棱 PA, PB, PC
\draw[black!70, thick] (P) -- (A);
\draw[black!70, thick] (P) -- (B);
\draw[black!70, thick] (P) -- (C);

% 中线 PE（红色）
\draw[red, thick] (P) -- (E);

% AE 辅助线（连接 A 到 BC 中点）
\draw[gray, dashed] (A) -- (E);

% 垂直符号（PE ⊥ 底面，在 E 处）
% 在 E 处画小方块，方向沿底面
\draw[gray, thin]
  ($(E) + (0.15, 0)$) -- ($(E) + (0.15, 0.15)$) -- ($(E) + (0, 0.15)$);

% 直角符号 AE ⊥ BC（可选，体现三垂线）
% E 是 BC 中点，正三角形中 AE ⊥ BC
\draw[orange!80, thin]
  ($(E) + (-0.13, 0.07)$) -- ($(E) + (-0.06, 0.20)$) -- ($(E) + (0.07, 0.13)$);

% 标记顶点
\node[dot] at (P) {};
\node[dot] at (A) {};
\node[dot] at (B) {};
\node[dot] at (C) {};
\node[dot] at (E) {};

% 标注
\node[font=\small, above] at (P) {$P$};
\node[font=\small, left] at (A) {$A$};
\node[font=\small, right] at (B) {$B$};
\node[font=\small, above right] at (C) {$C$};
\node[font=\small, below right] at (E) {$E$};

% 底部说明
\node[font=\small, align=center] at (0, -0.7)
  {正三棱锥 $P$-$ABC$，$E$ 为 $BC$ 中点};
\node[font=\small, align=center] at (0, -1.1)
  {$PE \perp$ 平面 $ABC$};

\end{tikzpicture}
\end{document}
